\documentclass[11pt]{article}
\usepackage[margin=1in]{geometry}
\usepackage{amsmath, amssymb, amsthm}
\usepackage{mathtools}

\newtheorem{theorem}{Theorem}
\newtheorem{lemma}[theorem]{Lemma}
\newtheorem{proposition}[theorem]{Proposition}
\theoremstyle{definition}
\newtheorem{definition}[theorem]{Definition}
\newtheorem{example}[theorem]{Example}
\theoremstyle{remark}
\newtheorem{remark}[theorem]{Remark}

\newcommand{\R}{\mathbb{R}}
\newcommand{\norm}[1]{\left\lVert #1 \right\rVert}
\newcommand{\grad}{\nabla}
\DeclareMathOperator{\dist}{dist}

\title{Multivariate Calculus: Partials, Total Derivatives, and the Chain Rule}
\author{Math Foundations --- Node 13}
\date{}

\begin{document}
\maketitle

\section{Definitions}

Throughout, $U \subset \R^n$ is open and $f: U \to \R$. We write $\norm{\cdot}$ for the Euclidean norm and $e_1, \dots, e_n$ for the standard basis.

\begin{definition}[Partial derivative]
Let $x \in U$. The \emph{$i$-th partial derivative} of $f$ at $x$ is
\[
\frac{\partial f}{\partial x_i}(x) := \lim_{h \to 0} \frac{f(x + h e_i) - f(x)}{h},
\]
provided the limit exists. Equivalently, holding all coordinates fixed except the $i$-th and viewing $f$ as a univariate function of $x_i$, the partial is its ordinary derivative.
\end{definition}

\begin{definition}[Gradient]
When all $n$ partials of $f$ at $x$ exist, the \emph{gradient} of $f$ at $x$ is
\[
\grad f(x) := \left( \frac{\partial f}{\partial x_1}(x), \dots, \frac{\partial f}{\partial x_n}(x) \right) \in \R^n.
\]
\end{definition}

\begin{definition}[Directional derivative]
For a unit vector $v \in \R^n$, the \emph{directional derivative} of $f$ at $x$ in direction $v$ is
\[
D_v f(x) := \lim_{h \to 0} \frac{f(x + h v) - f(x)}{h}.
\]
\end{definition}

\begin{definition}[Total / Fr\'echet derivative]
$f$ is \emph{differentiable at} $x \in U$ if there exists a linear map $L : \R^n \to \R$ such that
\[
\lim_{h \to 0} \frac{\norm{f(x+h) - f(x) - L h}}{\norm{h}} = 0.
\]
If $L$ exists it is unique; we write $Df(x) := L$. Identifying linear maps $\R^n \to \R$ with row vectors via the inner product, $L h = \grad f(x) \cdot h$.
\end{definition}

\begin{remark}
Existence of all partials at $x$ does \emph{not} imply differentiability, nor even continuity. A standard counterexample is $f(x,y) = xy/(x^2+y^2)$ on the punctured plane, extended by $f(0,0)=0$: both partials vanish at the origin, but along $y = x$ the function equals $1/2$, so it is discontinuous there.
\end{remark}

\section{From partials to the total derivative}

\begin{proposition}[Continuous partials imply differentiability]\label{prop:c1}
If $\partial f / \partial x_i$ exist in an open neighborhood of $x$ and are continuous at $x$, then $f$ is differentiable at $x$ with $Df(x) h = \grad f(x) \cdot h$.
\end{proposition}

\noindent A proof is in Rudin, Theorem~9.21. In ML we work almost exclusively with $C^1$ functions, so this proposition lets us identify differentiability with the gradient picture without further worry.

\section{The chain rule}

The central result of this node is the multivariate chain rule along a curve.

\begin{theorem}[Multivariate chain rule, scalar codomain]\label{thm:chain}
Let $U \subset \R^n$ be open, $f : U \to \R$ differentiable at $a \in U$, and $\gamma : I \to U$ a curve on an open interval $I \subset \R$ with $\gamma(t_0) = a$ and $\gamma$ differentiable at $t_0$. Then $f \circ \gamma : I \to \R$ is differentiable at $t_0$ and
\[
(f \circ \gamma)'(t_0) \;=\; \grad f(a) \cdot \gamma'(t_0) \;=\; \sum_{i=1}^{n} \frac{\partial f}{\partial x_i}(a) \, \gamma_i'(t_0).
\]
\end{theorem}

\begin{proof}
Set $L := Df(a)$, identified with $\grad f(a) \in \R^n$ via $L w = \grad f(a) \cdot w$. Differentiability of $f$ at $a$ means: for every $\varepsilon > 0$ there is $\delta > 0$ such that
\begin{equation}\label{eq:diffF}
\norm{w} < \delta \;\Longrightarrow\; \bigl| f(a + w) - f(a) - L w \bigr| \le \varepsilon \norm{w}.
\end{equation}
Equivalently, there is a function $\varphi : \R^n \to \R$ with $\varphi(w) \to 0$ as $w \to 0$ and
\begin{equation}\label{eq:diffF2}
f(a + w) - f(a) = L w + \varphi(w) \norm{w}, \qquad \varphi(0) := 0.
\end{equation}

Differentiability of $\gamma$ at $t_0$ similarly gives: there is a function $\psi : I - t_0 \to \R^n$ with $\psi(s) \to 0$ as $s \to 0$ and
\begin{equation}\label{eq:diffGamma}
\gamma(t_0 + s) - \gamma(t_0) = s \, \gamma'(t_0) + s \, \psi(s), \qquad \psi(0) := 0.
\end{equation}

Set $w(s) := \gamma(t_0 + s) - a$. By~\eqref{eq:diffGamma}, $w(s) = s(\gamma'(t_0) + \psi(s))$, so
\begin{equation}\label{eq:wbound}
\norm{w(s)} \le |s| \bigl( \norm{\gamma'(t_0)} + \norm{\psi(s)} \bigr).
\end{equation}
Pick $s$ small enough that $\gamma(t_0 + s) \in U$ and $\norm{w(s)} < \delta$ (possible since $\gamma$ is continuous at $t_0$). Substituting $w = w(s)$ into~\eqref{eq:diffF2}:
\[
f(\gamma(t_0+s)) - f(a) = L w(s) + \varphi(w(s)) \norm{w(s)}.
\]
By linearity of $L$ and~\eqref{eq:diffGamma},
\[
L w(s) = L\bigl( s \gamma'(t_0) + s \psi(s) \bigr) = s \, L \gamma'(t_0) + s \, L \psi(s).
\]
Therefore
\begin{equation}\label{eq:combine}
\frac{f(\gamma(t_0+s)) - f(\gamma(t_0))}{s} = L \gamma'(t_0) + L \psi(s) + \varphi(w(s)) \cdot \frac{\norm{w(s)}}{s}.
\end{equation}

We send $s \to 0$ term by term. The first term $L \gamma'(t_0) = \grad f(a) \cdot \gamma'(t_0)$ is the target. For the second, $L$ is continuous (every linear map on $\R^n$ is), and $\psi(s) \to 0$, so $L \psi(s) \to 0$. For the third, $w(s) \to 0$ by~\eqref{eq:wbound}, so $\varphi(w(s)) \to 0$, while $|\norm{w(s)}/s|$ is bounded above by $\norm{\gamma'(t_0)} + \norm{\psi(s)}$, which stays bounded near $s = 0$; hence the product tends to $0$.

Combining, the limit of the left-hand side of~\eqref{eq:combine} exists and equals $\grad f(a) \cdot \gamma'(t_0)$. That limit is by definition $(f \circ \gamma)'(t_0)$, completing the proof.
\end{proof}

\begin{remark}[The directional-derivative corollary]
Take $\gamma(t) = x + t v$ with $v \in \R^n$ a unit vector. Then $\gamma'(0) = v$, and Theorem~\ref{thm:chain} specializes to
\[
D_v f(x) = (f \circ \gamma)'(0) = \grad f(x) \cdot v.
\]
This identity is the basis of the Cauchy--Schwarz argument that the gradient points in the direction of steepest ascent: $D_v f(x) \le \norm{\grad f(x)}$ with equality at $v = \grad f(x)/\norm{\grad f(x)}$.
\end{remark}

\section{Worked example}

Let $f(x,y) = x^2 y + \sin(xy)$. Computing partials by holding the other variable fixed:
\[
\frac{\partial f}{\partial x} = 2xy + y\cos(xy), \qquad \frac{\partial f}{\partial y} = x^2 + x\cos(xy).
\]
At $(1,2)$: $\partial_x f = 4 + 2\cos 2$, $\partial_y f = 1 + \cos 2$. So $\grad f(1,2) = (4 + 2\cos 2,\; 1 + \cos 2)$. For $v = (1,1)/\sqrt 2$,
\[
D_v f(1,2) = \frac{(4 + 2\cos 2) + (1 + \cos 2)}{\sqrt 2} = \frac{5 + 3\cos 2}{\sqrt 2}.
\]
The companion script \texttt{code.py} confirms these values numerically.

\end{document}
